\documentclass[../main.tex]{subfiles}

\begin{document}

\section{Phased Mitigation Roadmap}
\label{app:phased-roadmap}

This appendix turns Table~\ref{tab:findings-response-routing} into a proposed
delivery plan; it does not show stakeholder acceptance or implementation.
Phase~0--2 clocks begin when an affected deployment or equivalent exposure is
identified. The \code{RM-1A} census, \code{RM-1E} assertion catalog, evidence
actions, and Phase~3 programme begin at roadmap adoption. Adoption must assign
dated owners and inventory deadlines; a pre-adoption exposure enters its
response route immediately.

Phase, dependency, owner, effort, and sequence are treatment-planning fields,
not risk inputs. A blocked final control does not extend a response clock: the
owner maintains tested containment, records the variance, and obtains an
explicit residual-risk decision. At publication every package is
\emph{Proposed}, implementation is \emph{not verified}, and no residual risk
has been accepted. Within ten business days of adoption, \code{GOV} should add
a named accountable person, programme record, date or release, and status to
each untriggered package. A triggered Phase~0 record receives its \code{APP}
owner and \code{GOV} authority at T0.

\subsection{Accountability Roles}

\begingroup
\small
\setlength{\tabcolsep}{3pt}
\renewcommand{\arraystretch}{1.1}
\setlength{\LTleft}{-0.7in}
\setlength{\LTright}{-0.7in}
\begin{longtable}{@{}p{0.12\widetablewidth}p{0.27\widetablewidth}p{0.55\widetablewidth}@{}}
\caption{Proposed roadmap accountability roles}
\label{tab:roadmap-roles}\\
\hline
\textbf{Code} & \textbf{Role} & \textbf{Accountability} \\
\hline
\endfirsthead
\caption[]{Roadmap roles, continued}\\
\hline
\textbf{Code} & \textbf{Role} & \textbf{Accountability} \\
\hline
\endhead
\hline
\endfoot
\hline
\endlastfoot
\code{GOV} & Roadmap and risk authority & Scope, assignments, gates,
exceptions, and residual-risk decisions. \\
\code{CORE} & Core library and tools & Parser, arithmetic, policy API, tool,
regression, fuzz, and durability changes. \\
\code{APP} & Application/service/storage operations & Deployment policy,
containment, rollout, capacity, checkpoint, recovery, and independent copies. \\
\code{PRIV} & Privacy and data governance & Classification, minimization,
disclosure, retention, encryption, and privacy acceptance. \\
\code{INTEROP} & Format and interoperability & Specification, reference and
independent implementations, fixtures, repair, and migration. \\
\code{EXT} & Extension ecosystem & Plugin, VOL, VFD, and filter policy,
manifests, trust, introspection, and documentation. \\
\code{REL} & Release and supply chain & Provenance, dependencies, SBOMs,
release gates, and configuration reporting. \\
\code{ASSURE} & Independent verifier & Reproduces acceptance evidence and
signs the verification used for closure. \\
\end{longtable}
\endgroup

Each package has one accountable owner, plus a responsible implementer,
\code{ASSURE} verifier, and release or deployment approver. Cross-project work
is instantiated once per governed scope. For High or Critical scope, the
verifier is not also owner, rollout approver, or risk acceptor. Only
\code{GOV} approves gates, exceptions, or residual-risk dispositions.

\subsection{Phase Gates and Clocks}

Phases are delivery waves, not severity levels; work may run in parallel, but
a later structural package never extends an earlier response clock. Each
finding-and-scope instance therefore has an SLA child record containing its
scope, T0, route, due and performed dates, interim action, status, and evidence,
separate from a shared package's structural due date.

\begingroup
\small
\setlength{\tabcolsep}{3pt}
\renewcommand{\arraystretch}{1.1}
\setlength{\LTleft}{-0.8in}
\setlength{\LTright}{-0.8in}
\begin{longtable}{@{}p{0.14\widetablewidth}p{0.20\widetablewidth}p{0.27\widetablewidth}p{0.33\widetablewidth}@{}}
\caption{Roadmap phases and exit gates}
\label{tab:roadmap-phases}\\
\hline
\textbf{Phase} & \textbf{Clock} & \textbf{Outcome} & \textbf{Exit gate} \\
\hline
\endfirsthead
\caption[]{Roadmap phases and exit gates, continued}\\
\hline
\textbf{Phase} & \textbf{Clock} & \textbf{Outcome} & \textbf{Exit gate} \\
\hline
\endhead
\hline
\endfoot
\hline
\endlastfoot
0---Contain & Before the next untrusted job and no later than 24 hours after
identification & Remove the external-resource path from attacker reach;
preserve evidence and rotate exposed secrets where indicated. & No confirmed
Critical-equivalent path remains active without tested containment; owner,
scope, and decision are recorded. \\
1---Verify/remediate & Census and assertion catalog within 30 days of adoption;
Near-term treatment within 30 days of affected-scope identification & Bound the
population and apply a supported fix or tested operational control. & Each
applicable finding has scope evidence, a negative test or rehearsed control,
implementation record, and residual-risk disposition; confirmed assertion
paths are routed without waiting for global migration. \\
2---Release safeguards & Planned treatment within 90 days of identification or
the next supported release/migration window, whichever is sooner; evidence
actions target 90 days after adoption & Deliver strict-processing,
interoperability, archive, and bounded evidence packages. & Tests pass;
compatibility cases have deterministic outcomes; evidence records are promoted
and routed, negatively dispositioned, or remain bounded and open. \\
3---Recurring classes & Target 3--12 months after adoption through staged
supported releases & Deliver invariant parsing, validation, fuzz/review gates,
durability architecture, and governed extensions. & Each named tranche passes
acceptance; untreated tranches remain visible with containment and dates. \\
4---Sustain/reassess & PR, nightly, release, and at least annual cadence as
applicable & Detect regressions and ecosystem mismatch; maintain verified
residual decisions. & Gates run; failures block or receive named, expiring
exceptions; registers are updated. \\
\end{longtable}
\endgroup

When a package covers several routes, each finding keeps its own maximum
horizon. Co-delivery does not raise or lower urgency, and a package cannot hide
an overdue SLA child behind its later architectural date.

\subsection{Licensing Constraint on Reuse of Existing Work}
\label{sec:roadmap-licensing}

\code{RM-1E}, \code{RM-2A}, \code{RM-2C}, and \code{RM-3A}--\code{RM-3C}
may draw on the H5Lens / \code{hdf5-pickles} decoder, invariant registry,
preflight oracle, corpus, and measurements
\cite{hdf5-pickles,h5policy,h5policy-invariant-registry}. H5Lens reports
GPLv3-or-later licensing while HDF5 uses a permissive BSD-style license. This
is an implementation, provenance, and governance constraint, not a legal
conclusion. Before reuse, \code{GOV}/\code{REL} must inventory rights and
licenses, record provenance per artifact, and obtain qualified review where
needed.

Planning assumes only that approved factual interfaces, reproducible fixtures,
and measurements may inform independent native implementation. It does not
authorize source transfer, linking, vendoring, redistribution, or relicensing.
Invariant vocabulary is candidate design input; an external oracle is evaluated
separately from component reuse; each fixture receives its own provenance and
redistribution decision; and case records remain triage inputs requiring local
reachability, version, and consequence verification. If review rejects these
uses, affected effort estimates are revised.

\subsection{Work-Package Register}

This is the authoritative package, dependency, owner, phase, and effort
register. Finding coverage is indexed in
Table~\ref{tab:roadmap-finding-coverage}; recommendation coverage is indexed in
Table~\ref{tab:recommendation-register}; observable acceptance is stated once
in Table~\ref{tab:roadmap-acceptance}. Dependencies enable final controls but
never defer containment. Effort is breadth, not risk or elapsed time.

\begingroup
\small
\setlength{\tabcolsep}{2pt}
\renewcommand{\arraystretch}{1.18}
\setlength{\LTleft}{-0.9in}
\setlength{\LTright}{-0.9in}
\begin{longtable}{@{}>{\raggedright\arraybackslash}p{0.07\widetablewidth}>{\raggedright\arraybackslash}p{0.11\widetablewidth}>{\raggedright\arraybackslash}p{0.46\widetablewidth}>{\raggedright\arraybackslash}p{0.25\widetablewidth}>{\raggedright\arraybackslash}p{0.07\widetablewidth}@{}}
\caption{Proposed mitigation work packages}
\label{tab:roadmap-work-packages}\\
\hline
\textbf{ID} & \textbf{Phase / owner} & \textbf{Objective} &
\textbf{Dependencies} & \textbf{Effort} \\
\hline
\endfirsthead
\caption[]{Mitigation work packages, continued}\\
\hline
\textbf{ID} & \textbf{Phase / owner} & \textbf{Objective} &
\textbf{Dependencies} & \textbf{Effort} \\
\hline
\endhead
\hline
\endfoot
\hline
\endlastfoot
\code{RM-0A} & 0 / \code{APP} & Disable or sandbox all supported external
resource forms in untrusted processing; remove unnecessary worker authority,
preserve evidence, and rotate indicated secrets. & None; isolate or stop if no
library control exists. & Medium \\
\code{RM-1A} & 1 / \code{GOV} & Census deployments, supported versions,
features, provenance, authority, scope tags, build/assertion mode, owners, and
current control evidence. & Deployment and project evidence. & Medium \\
\code{RM-1B} & 1 / \code{CORE} & Patch or isolate affected core/tools; replace
confirmed sole assertion guards; retain minimized debug/\code{-DNDEBUG} and
sanitizer regressions. & \code{RM-1A}; promoted \code{RM-1E} cases; supported
fix or containment. & Medium \\
\code{RM-1C} & 1 / \code{APP} & Patch or contain affected loaders; test safe
deserialization, external-reference denial, allocation budgets, and plugin-path
ownership. & Triggered \code{RM-0A}; \code{RM-1A}; supported release or local
control. & Medium \\
\code{RM-1D} & 1 / \code{PRIV} & Bound sensitive metadata and recipients;
minimize/tokenize or restrict the sharing boundary. & \code{RM-1A}; data owner
and classification. & Medium \\
\code{RM-1E} & 1 / \code{CORE} & Within 30 days of adoption, catalog supported
file-byte-dependent assertion sites, guards, reachability, sinks, paired build
results, versions, consequences, and dispositions; immediately route confirmed
adverse paths. & Supported-branch owners and reproducible builds. & Medium \\
\code{RM-1F} & 1 / \code{APP} & Within 30 days of deployment identification,
test headroom alerts, checkpoint/quiescence discipline, recovery rehearsals,
and an independent copy against stated recovery objectives. & Deployment,
storage, and data owners; census is not prerequisite for known scope. & Medium \\
\code{RM-2A} & 2 / \code{CORE} & Ship an enforceable untrusted-input profile
and identity-bound preflight with external-resource, plugin, traversal,
metadata, allocation, and expansion budgets; harden tools. & \code{RM-1A};
\code{RM-1B} fixtures are coordination input. & Large \\
\code{RM-2B} & 2 / \code{APP} & Integrate runtime denial and identity-bound
preflight; enforce overflow-safe object and aggregate budgets before
materialization. & \code{RM-1A}, \code{RM-1C}; coordinate with \code{RM-2A},
using an interim wrapper if needed. & Medium \\
\code{RM-2C} & 2 / \code{INTEROP} & Maintain valid, invalid, and contested
fixtures across supported readers/writers; inventory and disposition legacy
files; adjudicate each identified spec/library dispute within 90 days of
adoption. & \code{RM-1A}; external readers, writers, and fixtures. & Large \\
\code{RM-2D} & 2 / \code{CORE} & Inject storage exhaustion and abrupt
termination across write/flush/close, versions, VFDs, filesystems, and buffering;
measure errors, durable state, reopen/repair, read-back divergence, and loss. &
Fault injection, representative workloads, and data-owner input. & Large \\
\code{RM-2E} & 2 / \code{PRIV} & Trace one bounded extension deployment's data,
credentials, caches, logs, networks, recipients, and retention to a consequence
or defensible negative result. & Sponsor, \code{EXT}, classification, and
observable environment. & Medium \\
\code{RM-2F} & 2 / \code{REL} & Reproduce a producer/consumer mismatch with
exact profiles, artifacts, failure, recoverability, population, and recurrence.
& \code{INTEROP}, downstream owners, builds, and fixtures. & Medium \\
\code{RM-2G} & 2 / \code{PRIV} & Where minimization and recipient controls are
insufficient, deploy all-metadata protection with tested workflow, custody,
backup, and recovery. & \code{RM-1D}; supported design and key custodian. & Large \\
\code{RM-2H} & 2 / \code{INTEROP} & Within 90 days of archive identification,
bound the corpus, manifest every non-core dependency, test absence, classify
self-contained/external views, and disposition unavailable/removable
dependencies. & Repository owner and preservation objective; other package
inputs are non-blocking. & Medium \\
\code{RM-3A} & 3 / \code{CORE} & Deliver versioned invariants and mandatory
checked arithmetic, bounds, allocation, and traversal primitives, each mapped
to always-active enforcement. & \code{RM-1B}; \code{RM-1E}/\code{RM-2C}
coordination by tranche. & Large \\
\code{RM-3B} & 3 / \code{CORE} & Migrate prioritized structures to bounded raw
decode, semantic validation, then construction, in releasable tranches. &
\code{RM-3A}. & Large \\
\code{RM-3C} & 3 / \code{CORE} & Make class closure auditable through
structure-aware fuzzing, permanent regression, parser/tool review gates, build-
mode parity, and public metrics. & \code{RM-3A}; harness/corpus may start
earlier. & Large \\
\code{RM-3D} & 3 / \code{EXT} & Establish plugin identities, protected paths,
manifests, resource limits, signature tests, and bounded extension privacy
documentation/introspection. & \code{RM-2A}; privacy priority from \code{RM-2E}. & Large \\
\code{RM-3E} & 3 / \code{REL} & Publish SBOMs, pinned automation, signatures,
provenance, and feature reports; maintain a tested producer/consumer profile
matrix and deterministic mismatch dispositions. & \code{RM-1A}; mismatch
evidence from \code{RM-2F}. & Large \\
\code{RM-3F} & 3 / \code{PRIV} & Implement evidenced extension minimization,
temporary-artifact handling, retention/deletion, and protection of sensitive
filter metadata. & \code{RM-1D}, \code{RM-2E}; encryption uses \code{RM-2G}. & Large \\
\code{RM-3G} & 3 / \code{CORE} & Deliver VFD-level storage-exhaustion
notification, suspect-file state, optional reservation/pre-flush controls, a
per-layout/driver crash guarantee, and selected crash-proofing mechanism. &
\code{RM-2D}; a pass-through VFD may prove the design. & Large \\
\code{RM-3H} & 3 / \code{INTEROP} & Automate per-file dependency reports;
define an archival profile and deposit classification; preserve reviewed codec
specifications and test vectors. & \code{RM-2H}, \code{RM-2C}, \code{RM-3E},
and repository partner. & Medium \\
\code{RM-3I} & 3 / \code{INTEROP} & Maintain a declared machine-readable
format subset generating validators/fixtures and a standing adjudication
procedure; \code{GOV} decides whether it is normative. & \code{RM-2C}; reuse
requires Section~\ref{sec:roadmap-licensing} review. & Large \\
\code{RM-4A} & 4 / \code{GOV} & Operate gates and metrics; rerun applicability
and residual reviews; sustain corpora and close or re-plan exceptions. & Begins
with Phase~1 and absorbs accepted packages. & Medium recurring \\
\end{longtable}
\endgroup

\subsection{Minimum Viable Subset Under Constrained Capacity}
\label{sec:roadmap-minimum-viable}

Capacity ordering is not risk ranking and never changes a route or clock.
Phase~0--1 obligations and the bounded Planned deliverables remain due even if
the structural programme is deferred. The prevention ordering below is the
assessor's qualitative judgment under the stated reuse assumptions; no
engineer-week, cost, or comparative-effectiveness model was available.

\begingroup
\small
\setlength{\tabcolsep}{3pt}
\renewcommand{\arraystretch}{1.1}
\setlength{\LTleft}{-0.8in}
\setlength{\LTright}{-0.8in}
\begin{longtable}{@{}p{0.15\widetablewidth}p{0.40\widetablewidth}p{0.39\widetablewidth}@{}}
\caption{Capacity-constrained programme ordering}
\label{tab:roadmap-minimum-viable}\\
\hline
\textbf{Tier} & \textbf{Items and basis} & \textbf{Boundary if deferred} \\
\hline
\endfirsthead
\caption[]{Capacity-constrained ordering, continued}\\
\hline
\textbf{Tier} & \textbf{Items and basis} & \textbf{Boundary if deferred} \\
\hline
\endhead
\hline
\endfoot
\hline
\endlastfoot
A---not discretionary & Triggered \code{RM-0A}; \code{RM-1A}--\code{RM-1F}.
These contain Critical-equivalent paths, establish population, treat known
scope, contain write-side hazards, and bound assertions. & Omission breaches an
Immediate or Near-term clock. Census incompleteness cannot be used to defer a
known deployment. \\
Bounded Phase~2 & Case-specific \code{RM-2C} treatment and \code{RM-2H}
archive disposition complete within their Planned 90-day clocks. & Broader
matrices, machine-readable format work, automated reporting, and archival
profile design may be sequenced; the bounded deliverables may not. \\
B---fund first & \code{CORE-S6}, \code{CORE-S7}, \code{CORE-S3},
\code{CORE-S4}, then \code{CORE-S1} through their mapped packages. Existing
corpus, profile, and invariant artifacts may reduce design/enumeration work
after review; checked primitives are tranchable. & Regression, change-time
gates, shared policy, and measurable invariant coverage remain incomplete;
local treatments stay open as local rather than class closure. \\
C---explicit deferral & Remaining structural packages, including parser
migration, expanded fuzzing, plugin/provenance, encryption, full differential
and profile matrices, durability architecture, extension privacy, archive
automation, and machine-readable format governance. & Record a scoped,
expiring exception and retain interim controls. Deferral prevents the
corresponding class, lifecycle, durability, or governance closure claim. \\
\end{longtable}
\endgroup

Funding A, bounded Phase~2, and B is coherent but does not eliminate recurring
classes. \code{CORE-S0} prevention records remain open and \code{CORE-S12}
metrics must show the untreated denominator.

\subsection{Acceptance and Closure Evidence}

Table~\ref{tab:roadmap-acceptance} is the canonical acceptance specification.
For every row, the implementer links artifacts, \code{ASSURE} reproduces or
independently verifies them, the package owner attests readiness, and the
release/deployment approver authorizes rollout. \code{GOV} then records the
credited scope and a current residual-risk decision. Guidance or code shipment
alone is insufficient; failed, incomplete, or inconclusive tests remain open.

\begingroup
\normalsize
\setlength{\tabcolsep}{3pt}
\renewcommand{\arraystretch}{1.18}
\setlength{\LTleft}{-0.8in}
\setlength{\LTright}{-0.8in}
\begin{longtable}{@{}p{0.11\widetablewidth}p{0.83\widetablewidth}@{}}
\caption{Required work-package acceptance evidence}
\label{tab:roadmap-acceptance}\\
\hline
\textbf{Package} & \textbf{Required observable evidence} \\
\hline
\endfirsthead
\caption[]{Work-package acceptance evidence, continued}\\
\hline
\textbf{Package} & \textbf{Required observable evidence} \\
\hline
\endhead
\hline
\endfoot
\hline
\endlastfoot
\code{RM-0A} & Configuration/code diff and deployment proof; negative fixtures
for every supported external-resource form; sandbox-escape test; indicated
credential/secret rotation. \\
\code{RM-1A} & Machine-readable inventory of version, linkage, features,
build/assertion mode, scope tag, trust/authority, exposure, owner, boundary,
completeness attestation, and evidence date. \\
\code{RM-1B} & Version matrix, patch/isolation record, minimized fixtures, and
assertion-enabled plus \code{-DNDEBUG} ASan/UBSan/LSan controlled-error results. \\
\code{RM-1C} & Rollout evidence and negative tests for executable reconstruction,
external paths, allocation budgets, and plugin paths; startup state and canaries
in CWD, environment-selected, and writable searched paths. \\
\code{RM-1D} & Classified metadata/data-flow and recipient matrix; before/after
HDF5-API and raw-byte inspection; decision whether \code{RM-2G} is required. \\
\code{RM-1E} & Versioned site catalog with branch, function, predicate,
file-byte provenance, dominating guard, first sink, reachability, paired build
results, version scope, consequence, and disposition; tests and promoted
finding/SLA children for confirmed adverse paths. \\
\code{RM-1F} & Defined scope and recovery objectives; alert before injected
exhaustion; checkpoint/quiescence runbook; successful exhaustion and abrupt-
termination rehearsals; verified independent-copy recency, integrity,
location, and failure-domain separation. \\
\code{RM-2A} & API/policy documentation, default-state tests, hostile corpus,
resource/plugin/external-reference tests, file-replacement identity test, and
tool-default tests. \\
\code{RM-2B} & Per-downstream integration and deployment evidence; extreme,
overflow, cumulative, and valid-boundary allocation tests before materialization;
external-reference and file-replacement tests. \\
\code{RM-2C} & Versioned fixture metadata and reader matrix; exact release/
backport verification; \#6409 sufficient/undersized pair; corrected-writer
round trip; complete legacy enumeration and checksum/provenance-preserving
repair or disposition; dated rulings naming authority and changed artifact. \\
\code{RM-2D} & Fault-injection matrix for both triggers across API sequence,
version, VFD/filesystem/buffering profile, recording surfaced error, durable
state, reopen/repair, acknowledged-versus-read-back divergence, uniqueness,
and workload boundary. \\
\code{RM-2E} & Reproducible extension data-flow test with classified population,
recipients, endpoints, caches/logs/artifacts, retention, controls, and bounded
consequence or documented negative result. \\
\code{RM-2F} & Exact producer/consumer profiles, reproducible fixture and
diagnostic, artifact inventory, recoverability, population, recurrence, and
documented negative or promoted result. \\
\code{RM-2G} & Ciphertext/metadata inspection, unauthorized-open and partial-I/O
tests, protected key inventory, backup/restore, and successful key-loss recovery. \\
\code{RM-2H} & Defined corpus and completeness attestation; per-file manifest;
self-contained/external-view classification; absent-dependency tests; dated
owner disposition; provenance, license-review, integrity, and reconstruction
evidence for retained runtime artifacts. \\
\code{RM-3A} & Versioned invariants, code mapping, checked primitives, migration
inventory, single-invariant rejection in both build modes, and proof that no
declared invariant relies only on an assertion. \\
\code{RM-3B} & Per-component decode/validate/construct boundary; no native
object before validation; valid/malformed corpus plus build-mode sanitizer and
fuzz evidence. Closure names migrated components only. \\
\code{RM-3C} & PR/nightly/weekly/release/CVE-closure jobs, minimized seeds,
structure coverage, review records, metrics, build-mode parity, and a test
showing an assertion-only file predicate fails the gate. \\
\code{RM-3D} & Plugin manifest/policy, protected allowlist or keystore, rejection
of unsigned/untrusted/writable-path/over-budget/wrong-profile plugins, and
extension documentation/introspection tests. \\
\code{RM-3E} & SBOMs, checksums, signatures, verified provenance, pinned
automation, dependency scans, feature report, finite profile catalog, fixtures,
CI results, and deterministic disposition for each unsupported pair. \\
\code{RM-3F} & Tests showing minimized extension metadata/caches, secure
temporary handling, enforced retention/deletion, and, where needed, ciphertext
inspection plus key recovery. \\
\code{RM-3G} & Distinguishable storage-exhaustion errors per supported VFD/
buffering profile; tested suspect-state lifecycle; published layout/driver
guarantee; writer-kill outcomes and read-back divergence. Reported success
followed by inconsistency blocks acceptance. \\
\code{RM-3H} & Generated per-file dependency reports verified by absence tests;
archival-profile files readable with non-core plugins disabled; deposit
classifications; independent codec reproduction from archived specification
and vectors. \\
\code{RM-3I} & Declared machine-readable subset and generated fixtures;
divergences classified as library defect, description/specification defect, or
explicit implementation-defined behavior with documented consequence;
adjudication procedure, decision on normative status, and current dispute
ledger. An unruled dispute remains open. \\
\code{RM-4A} & Production-equivalent cadenced results, trends, updated mappings,
expired-exception report, and signed point-in-time residual decisions. \\
\end{longtable}
\endgroup

\subsection{Finding and Evidence-Action Coverage}

A formal finding-and-scope child closes only when the affected population is
treated or shown inapplicable, its applicable package evidence passes, and
\code{GOV} records current residual risk. This never changes the historical
inherent rating. Structural work stays open independently unless explicitly
credited. An evidence record closes only with a bounded positive, negative, or
promoted disposition; inconclusive work is re-planned. Table entries are ID-only
indexes into the authoritative package and acceptance tables.

\begingroup
\normalsize
\setlength{\tabcolsep}{3pt}
\renewcommand{\arraystretch}{1.18}
\setlength{\LTleft}{-0.8in}
\setlength{\LTright}{-0.8in}
\begin{longtable}{@{}p{0.20\widetablewidth}p{0.36\widetablewidth}p{0.38\widetablewidth}@{}}
\caption{Register-to-roadmap coverage}
\label{tab:roadmap-finding-coverage}\\
\hline
\textbf{Record} & \textbf{Bounded treatment or evidence} &
\textbf{Continuing prevention / boundary} \\
\hline
\endfirsthead
\caption[]{Register-to-roadmap coverage, continued}\\
\hline
\textbf{Record} & \textbf{Bounded treatment or evidence} &
\textbf{Continuing prevention / boundary} \\
\hline
\endhead
\hline
\endfoot
\hline
\endlastfoot
\code{CORE-SEC-001} & \code{RM-1A}, \code{RM-1B} &
\code{RM-2A}, \code{RM-3A}--\code{RM-3C}, \code{RM-4A}; only named tranches close. \\
\code{CORE-SEC-002} & \code{RM-1A}, \code{RM-1B} &
\code{RM-2A}, \code{RM-2C}, \code{RM-3A}--\code{RM-3C}, \code{RM-4A}. \\
\code{CORE-TOOL-001} & \code{RM-1A}, \code{RM-1B} &
\code{RM-2A}, \code{RM-3C}, \code{RM-4A}. \\
\code{CORE-PRIV-001} & \code{RM-1A}, \code{RM-1D}; conditional \code{RM-2G} &
\code{RM-4A}; closure states the tested disclosure boundary. \\
\code{CORE-SAFE-001} & \code{RM-1A}, \code{RM-1F}, \code{RM-2D} &
\code{RM-3G}; closure applies only to the named deployment and measured profile. \\
\code{CORE-SAFE-002} & \code{RM-1A}, \code{RM-1F}, \code{RM-2D} &
\code{RM-3G}; structural closure names the delivered mechanism/layout/driver. \\
\code{LONG-SAF-001} & \code{RM-2H} &
\code{RM-3H}, \code{RM-2C}, \code{RM-3E}, \code{RM-4A}; closure is corpus-specific. \\
\code{APP-SEC-001} & \code{RM-0A}, \code{RM-1A}, \code{RM-1C} &
\code{RM-2A}, \code{RM-2B}, \code{RM-4A}; each governed service closes separately. \\
\code{APP-SEC-002} & \code{RM-1A}, \code{RM-1C} & \code{RM-4A}. \\
\code{APP-SEC-003} & \code{RM-1A}, \code{RM-1C}; \code{RM-0A} for a confirmed
Critical-equivalent service & \code{RM-2A}, \code{RM-2B}, \code{RM-4A}. \\
\code{APP-SEC-004} & \code{RM-1A}, \code{RM-1C} &
\code{RM-2A}, \code{RM-2B}, \code{RM-4A}. \\
\code{APP-SEC-005} & \code{RM-1A}, \code{RM-1C} & \code{RM-3D}, \code{RM-4A}. \\
\code{IND-SAF-001} & \code{RM-1A}, \code{RM-2C} & \code{RM-3A}, \code{RM-4A}. \\
\code{IND-SAF-002} & \code{RM-1A}, \code{RM-2C} & \code{RM-3A}, \code{RM-4A}. \\
\code{CORE-SEC-003} & Evidence \code{RM-1A}, \code{RM-1E}; conditional
\code{RM-1B} & \code{RM-3A}--\code{RM-3C}, \code{RM-4A}; each supported
adverse path is promoted and routed. \\
\code{FMT-SAF-001} & Administrative adjudication \code{RM-2C} &
\code{RM-3I}, \code{RM-4A}; closure is dispute/subset-specific and assigns no urgency. \\
\code{EXT-PRIV-001} & Evidence \code{RM-2E} & Conditional \code{RM-3D},
\code{RM-3F}; rate only after population and harm are supported. \\
\code{SUP-SAFE-001} & Evidence \code{RM-1A}, \code{RM-2F} & Conditional
\code{RM-3E}; rate only after a bounded adverse scenario is supported. \\
\end{longtable}
\endgroup

\subsection{Change Control and Residual Risk}
\label{sec:roadmap-change-control}

Each implementation record contains package, accountable owner, implementer,
verifier, rollout approver, scope, linked records and SLA children, trigger and
T0, route and due date, dependencies and owners, effort assumptions, status,
evidence, target and verified residual tier/status, risk authority, exception,
reassessment trigger, and decision date. Status is \emph{Proposed},
\emph{Accepted}, \emph{In progress}, \emph{Blocked}, \emph{Implemented pending
verification}, \emph{Verified}, \emph{Deferred under exception}, \emph{Risk
accepted}, or \emph{Superseded}. A missed date remains immutably
\emph{Overdue}; later exception does not erase it.

Every exception names \code{GOV}, rationale, exact scope and unmet requirement,
tested interim controls, monitoring owner/cadence, issue and expiry dates, and
automatic reopen triggers. It neither revises inherent risk nor permits a
known uncontained Critical-equivalent path to keep processing untrusted input.
Exceptions are finite.

After verification and rollout, \code{GOV} reassesses current residual risk
under Section~\ref{sec:risk-rubric}, records credited controls and a decision to
accept, reduce, avoid, transfer, or escalate, and preserves the inherent rating.
The result is a dated immutable snapshot, not a permanent closure claim.
Phases~0--2 are reviewed at least weekly while open, Phase~3 at release gates,
and Phase~4 at its cadence. Control failure, new incident/reproducer, changed
release/default/trust/privilege/automation/sensitivity/dependency, test or
telemetry regression, and exception expiry trigger reassessment and a new SLA
child while preserving history.

\end{document}
